\section{Análisis de casos y siguientes pasos}

\subsection{Arc AGI Puzzle Walkthrough}
En esta curva de benchmark de Arc AGI, podemos usar los siguientes pasos para impulsar el inference loop: 1) en el prompt se declara de antemano que este puzzle es high-difficulty; 2) en la etapa de sampling se generan entre 30 y 40 candidate response; 3) en la etapa de search estos candidate se toman como tree seed, y con un stepwise scorer se calcula el promisingness además de introducir un probability threshold; 4) en la etapa de iteration se ejecuta trace feedback, se autorrevisa la reasoning chain de cada response, y con ayuda del verifier margin se orienta la elección por votación mayoritaria. Todo el proceso también registra every candidate's token cost y latency, para future tuning.

\begin{importantbox}{Observación del pipeline de Arc AGI}
Hay que considerar Arc AGI como un campo de práctica de todo el proceso: prompt → sampling → search → trace → voting, cada eslabón necesita instrumentation; si el token cost se duplica de pronto, hay que revisar si el sampling diversity y el verifier margin se han derrumbado, y luego usar el runbook para volver al conservative mode.
\end{importantbox}

\subsection{Despliegue en tiempo de inferencia para tareas de alto valor}
Para tareas competitivas de Code/Math como AlphaProof, el pipeline normalmente también agrega extra verification: 1) Run deterministic verifier (e.g. proof checker) as soon as the candidate completes; 2) If verification fails, reroute to alternative candidate with similar margin; 3) Log the mismatch and feed trace back to sampling controller for prompt iteration. Este tipo de flujo convierte el inference-time compute en un closed-loop control rather than a one-shot guess.

\begin{knowledgebox}{Puntos clave del despliegue de alto valor}
1. Usar el deterministic verifier como fallback, para evitar que el majority vote tome un wrong candidate como respuesta correcta; 2. usar trace replay para reproducir cada verification failure, de modo que se puedan generar nuevos examples mediante prompting; 3. registrar las metrics (token cost, verifier success, trace iteration) del close-loop pipeline en el performance dashboard.
\end{knowledgebox}

\subsection{Resumen de la sección}
El análisis de casos mostró cómo llevar a la práctica el compute pipeline en tiempo de inferencia en tareas del tipo Arc AGI y AlphaProof: primero se determina el nivel del prompt, luego se hace que sampling/search/trace constituyan los controles, y finalmente se escribe la retroalimentación de verification en el iteration loop para la mejora continua.

